\documentclass[10pt,a4paper]{amsart}
\usepackage[margin=19mm]{geometry}
\usepackage{amsmath,amssymb,mathtools}
\usepackage[scr=rsfs,cal=euler]{mathalfa}
\usepackage{xfrac}
\usepackage[unicode,hidelinks,bookmarks=false]{hyperref}
\newcommand{\ie}{i.e.\ }
\newcommand{\cf}{cf.\ }
\newcommand{\resp}{respectively\ }
\newcommand{\Subsection}{\subsection}
\providecommand{\nobreakdash}{\nobreak\hbox{-}}
\setlength{\emergencystretch}{2em}
\multlinegap0pt
\newcommand{\Gtwo}{\mathrm{G}_2}
\usepackage{amssymb}
\usepackage{enumerate}
\usepackage{float}
\usepackage{tikz}
\newtheorem{lem}{Lemma}
\newtheorem{theorem}{Theorem}
\newcommand{\R}{\mathbb R}
\renewcommand{\b}{\beta}
\newcommand{\la}{\langle}
\newcommand{\ra}{\rangle}
\title{Compact holonomy $\Gtwo$ manifolds need not be formal}
\author{Luc\'ia Mart\'in-Merch\'an}
\date{Source: arXiv:2409.04362v1 (CC BY 4.0)}
\begin{document}
\maketitle

\noindent\emph{Source and licence.} Compact holonomy $\Gtwo$ manifolds need not be formal. Version arXiv:2409.04362v1 (CC BY 4.0), distributed by arXiv under the Creative Commons Attribution 4.0 International licence, http://creativecommons.org/licenses/by/4.0/, as recorded in the arXiv licence metadata for this exact version and shown on the abstract page https://arxiv.org/abs/2409.04362v1. Full licence notice: \texttt{licenses/arxiv-2409.04362-CC-BY-4.0.txt}.\par\smallskip

\noindent\textit{Editorial scope.} Counted unit: the article's theorem theo:non-formal (source0.tex lines 989--992). The article's complete original proof of it is delivered with it (source0.tex lines 993--1035, one \texttt{proof} environment of 43 lines). No mathematics is written, repaired or re-derived: the statement and the proof are the article's own lines, in the article's own notation, and no retained line is substituted or re-flowed.  Retained uncounted as the source-local context the proof needs: lines 588--590; lines 717--720; lines 729--735; lines 761--763; lines 961--969.  Nothing was omitted from any retained span, so the package carries no omission record.  No retained line contains a citation, so the wrapper carries no bibliography and no bibliography entry.  No retained line contains a figure environment, an image-inclusion command or drawing code, so no image is delivered and the package declares no asset dependency.  Editorial formatting only, in addition: an AMS \texttt{amsart} preamble that restates the article's own declarations and macros used by the retained lines, namely the environments \texttt{lem}, \texttt{theorem} on separate counters, together with the standard packages those lines need.  The retained environments print in this document as Lemma 1, Theorem 1 in order of appearance, whereas the article prints the same items with the numbering of its own full document.  The complete native source of the article as distributed in the arXiv e-print for this exact version is bundled unchanged as \texttt{sources/164882/source0.tex}.
\begin{equation} \label{eqn:cohom}
H^k(\widetilde{M})\cong H^k(X)\oplus \oplus_{j=1}^{10} H^{k-2}(N_j)\otimes \la \mathbf{x}_j\ra.
\end{equation}

\begin{align}
\mathbf{x}_j^2=&-2\mathrm{Th}[N_j], \qquad \qquad \qquad \qquad  \mathbf{x}_{j}\cdot \mathbf{x}_{k}=0 \mbox{ if } j\neq k, \label{eqn:prod-1} \\
[\alpha] \cdot \mathbf{x}_{j}=& [\alpha|_{N_j}]\otimes \mathbf{x}_j, \mbox{ if } [\alpha] \in H^*(X). \label{eqn:prod-2}
\end{align}

\begin{lem} \label{lem:PDuals}
Denote $\b_1=i dt\wedge (dz_{123} - dz_{\bar{1}\bar{2}\bar{3}})$, and  $\b_2=idt \wedge (dz_{12\bar{3}} - dz_{\bar{1}\bar{2}3})$. Then 
$$
PD[N_1]=PD[N_2]=\beta_1+\beta_2, \qquad 
PD[N_3]=PD[N_j]=(\beta_1 - \beta_2)/2, \qquad 4 \leq j \leq 10.
$$
\end{lem}

\begin{equation} \label{eqn:closed-Massey}
    y= \xi_{12} \wedge \alpha_3 - \hat{\alpha}_1 \wedge  {\xi}_{23}
\end{equation}

\begin{lem} \label{lem:prim-wtM}
There is a form $\widetilde{\alpha} \in \Omega^3(\widetilde{M})$ such that $d(\widetilde{\alpha})=\tau_7^2-\tau_3^2$.
In particular, it is closed on $U_1\cup U_2$. In addition,
$$
\widetilde{\alpha} \wedge \tau_1= -4\mathrm{pr}^*_1(\varphi|_{N_1}) \wedge \tau_1, \qquad
\widetilde{\alpha} \wedge \tau_2= 4\mathrm{pr}^*_2(\varphi|_{N_2}) \wedge \tau_2,
$$
where $\mathrm{pr}_j=\pi_j \circ \rho \colon \rho^{-1}(V_j') \to N_j$.
\end{lem}

\begin{theorem} \label {theo:non-formal}
The triple Massey product $\la [\tau_1 + \tau_2], [\tau_7+ \tau_3], [\tau_7-\tau_3] \ra$ is not trivial.
Therefore, $\widetilde{M}$ is non-formal.
\end{theorem}

\begin{proof}
Note that $(\tau_1 + \tau_2)\wedge (\tau_7+ \tau_3) =0$ and $(\tau_7+ \tau_3)\wedge (\tau_7-\tau_3)=\tau_7^2-\tau_3^2=d\widetilde{\alpha}$. The cohomology class determined by expression \eqref{eqn:closed-Massey} is
$$
[y]= [\widetilde{\alpha}\wedge (\tau_1 + \tau_2)],
$$
and the triple Massey product is trivial if and only if $y\in \mathcal{I}$, where
$\mathcal{I}$ is the ideal generated by $\tau_1 + \tau_2$ and $\tau_7-\tau_3$. We prove by direct computation that $y\notin \mathcal{I}$. 

First, by Lemma \ref{lem:prim-wtM} we know
$$
[y]= -4[\mathrm{pr}_1^*(\varphi|_{N_1}) \wedge \tau_1] + 4 [\mathrm{pr}_2^*(\varphi|_{N_2}) \wedge \tau_2].
$$
Under the isomorphism in equation \eqref{eqn:cohom}, this class is $y'=-4[\varphi|_{N_1}] \otimes \mathbf{x}_1 + 4[\varphi|_{N_2}]\otimes \mathbf{x}_2$.
We now compute the image  $\mathcal{I}'$ of the degree-5 part of the ideal $\mathcal{I}$ under the isomorphism,
namely
$$
H^3(X)\cdot \la \mathbf{x}_1+\mathbf{x}_2 \ra \oplus
H^3(X) \cdot \la \mathbf{x}_7-\mathbf{x}_3 \ra
\oplus (\la dy_3 \otimes \mathbf{x}_1, dy_3 \otimes \mathbf{x}_2, dx_3 \otimes \mathbf{x}_3, \dots, dx_3 \otimes \mathbf{x}_{10} \ra)\cdot (\la \mathbf{x}_1+\mathbf{x}_2 ,\mathbf{x}_7-\mathbf{x}_3 \ra) .
$$

The proof of Lemma \ref{lem:PDuals} ensures that 
 $\int_{N_1} \xi= \int_{N_2} \xi$ and  $\int_{N_3} \xi= \int_{N_7} \xi$ for every $[\xi] \in H^3(X)$.
 Fix $[\xi] \in H^3(X)$ and set $\lambda_1,\lambda_3\in \R$ so that $[\xi|_{N_j}]=\lambda_j [\varphi|_{N_j}]$ with $j=1,3$. Then, 
 $\int_{N_2}\xi=\int_{N_1}\xi= \lambda_1 \int_{N_1} \varphi= \lambda_1 \int_{N_2} \varphi$ and similarly $\int_{N_7}\xi= \lambda_3 \int_{N_7} \varphi$, that is, $[\xi|_{N_2}]=\lambda_1 [\varphi|_{N_2}]$, and  $[\xi|_{N_7}]=\lambda_3 [\varphi|_{N_7}]$. Hence, using equation \eqref{eqn:prod-2} we conclude, 
$$
[\xi] \cdot (\mathbf{x}_1+\mathbf{x}_2)= \lambda_1 ([\varphi|_{N_1}]\otimes \mathbf{x}_1+[\varphi|_{N_2}]\otimes \mathbf{x}_2), \qquad 
[\xi] \cdot (\mathbf{x}_7-\mathbf{x}_3)= \lambda_3 ( [\varphi|_{N_7}]\otimes \mathbf{x}_7-[\varphi|_{N_3}]\otimes \mathbf{x}_3).
$$
Thus,
$
H^3(X)\cdot \la \mathbf{x}_1+\mathbf{x}_2 \ra= \la [\varphi|_{N_1}]\otimes \mathbf{x}_1+[\varphi|_{N_2}]\otimes \mathbf{x}_2  \ra$, and  
$H^3(X)\cdot \la \mathbf{x}_7-\mathbf{x}_3 \ra= \la [\varphi|_{N_7}]\otimes \mathbf{x}_7-[\varphi|_{N_3}]\otimes \mathbf{x}_3  \ra .
$

In addition, 
$(\la dy_3 \otimes \mathbf{x}_1, dy_3 \otimes \mathbf{x}_2, dx_3 \otimes \mathbf{x}_3, \dots, dx_3 \otimes \mathbf{x}_{10} \ra)\cdot (\la \mathbf{x}_1+\mathbf{x}_2 ,\mathbf{x}_7-\mathbf{x}_3\ra)  \subset H^5(X)$ because
 $\mathbf{x}_j\mathbf{x}_k \in H^4(X)$ as a consequence of equation \eqref{eqn:prod-1}. For instance, $(dy_3\otimes \mathbf{x}_1) \cdot \mathbf{x}_1$ is represented by $\mathrm{pr}_1^*(dx_3) \wedge \tau_1^2$, and it induces a cohomology class on $X$ by pushforward as it is supported around $O_1$, and vanishes on a neighborhood of $E_1$. However, 
 $$
 y'=-4[\varphi|_{N_1}] \otimes \mathbf{x}_1 +4[\varphi|_{N_2}]\otimes \mathbf{x}_2 \notin H^5(X) \oplus \la [\varphi|_{N_1}]\otimes \mathbf{x}_1+[\varphi|_{N_2}]\otimes \mathbf{x}_2, [\varphi|_{N_7}]\otimes \mathbf{x}_7-[\varphi|_{N_3}]\otimes \mathbf{x}_3 \ra \supset \mathcal{I}',
 $$
 so $[y]\notin \mathcal{I}$.
\end{proof}

\end{document}
